\chapter{Classical minimal surfaces by calculation}\label{sec:examples-minimal}
We calculate metrics, curvature, topology, and a stability distinction for explicit surfaces. The source is \cite{ChMin}, Sections 13--15, pp.~35--45. Our parametrizations are fixed below, so signs and scale factors can be checked independently of source coordinates. Minimality alone will not be used to assert embeddedness, completeness, or minimizing behavior.

\begin{lemma}[A conformal harmonic immersion is minimal]\label{lem:example-conformal-harmonic}
Let $X:U\subset\mathbb R^2\to\mathbb R^N$ be a smooth immersion with $|X_u|^2=|X_v|^2=\lambda^2>0$ and $\langle X_u,X_v\rangle=0$. Then
\[
 \vec H=\lambda^{-2}(X_{uu}+X_{vv}).
\]
Consequently $X$ is minimal exactly when all its coordinate functions are harmonic.
\end{lemma}
\begin{proof}
The induced metric is $g=\lambda^2(du^2+dv^2)$. Its Laplace--Beltrami operator is $\Delta_g=\lambda^{-2}(\partial_u^2+\partial_v^2)$, by the coordinate divergence formula. For an immersion, tracing the normal second derivative gives $\Delta_gX=\vec H$: the tangential Christoffel terms cancel when taking the covariant Hessian. Substitution proves the formula. This is \cite{ChMin}, Lemma 15.1. Our mean curvature vector is the trace, with no factor of one half.
\end{proof}

\begin{example}[The catenoid: an embedded complete minimal annulus]\label{ex:explicit-catenoid}
For $a>0$, define
\[
 C_a(\theta,v)=a(\cosh v\cos\theta,\cosh v\sin\theta,v),
 \qquad(\theta,v)\in(\mathbb R/2\pi\mathbb Z)\times\mathbb R.
\]
This is a proper embedded complete minimal surface. Its induced metric is
\[
 g=a^2\cosh^2v\,(d\theta^2+dv^2).
\]
The waist is the circle $v=0$ of radius $a$. Its two open halves are graphs $z=\pm a\operatorname{arcosh}(r/a)$ over $r>a$.
\end{example}
\begin{proof}
Differentiation gives orthogonal coordinate vectors of equal squared length $a^2\cosh^2v$. The sum of second derivatives is zero, so Lemma \ref{lem:example-conformal-harmonic} gives minimality. The height determines $v$, and the horizontal direction determines $\theta$ modulo $2\pi$; thus the map is injective. A bounded ambient set bounds $v$, and the angular factor is compact, proving properness. The metric dominates the complete cylinder metric $a^2(d\theta^2+dv^2)$; an intrinsic Cauchy sequence therefore converges in that cylinder, and local smooth equivalence gives convergence in $g$. Finally $r=a\cosh v$ gives the graph formulas. See \cite{ChMin}, Section 14.1 and Example 15.4.
\end{proof}

\begin{lemma}[Gauss curvature in a conformal chart]\label{lem:example-conformal-curvature}
For a smooth surface metric $g=e^{2h}(du^2+dv^2)$,
\[
 K=-e^{-2h}(h_{uu}+h_{vv}).
\]
Equivalently, for $\lambda=e^h$, $K=-\lambda^{-2}\Delta\log\lambda$, with the flat coordinate Laplacian on the right.
\end{lemma}
\begin{proof}
Use the orthonormal coframe $\theta^1=e^hdu$, $\theta^2=e^hdv$. The connection form $\omega=h_vdu-h_udv$ satisfies $d\theta^1=-\omega\wedge\theta^2$ and $d\theta^2=\omega\wedge\theta^1$. The curvature equation $d\omega=K\theta^1\wedge\theta^2$ now gives $-(h_{uu}+h_{vv})du\wedge dv=Ke^{2h}du\wedge dv$. This fixes the sign using the usual convention that a round sphere has positive Gauss curvature.
\end{proof}

\begin{proposition}[Curvature and total curvature of the catenoid]\label{prop:catenoid-curvature-example}
For $C_a$,
\[
 K=-\frac1{a^2\cosh^4v},\qquad |B|^2=\frac2{a^2\cosh^4v},
 \qquad\int_{C_a}K\,dA=-4\pi.
\]
The integral counts the catenoid once, with $0\leq\theta<2\pi$.
\end{proposition}
\begin{proof}
Apply Lemma \ref{lem:example-conformal-curvature} with $\lambda=a\cosh v$ and $(\log\lambda)_{vv}=\operatorname{sech}^2v$, giving $K$. Since the surface is minimal in $\mathbb R^3$, its principal curvatures are opposite and the Gauss equation gives $|B|^2=-2K$. Multiplying $K$ by $dA=a^2\cosh^2v\,d\theta\,dv$ leaves $-\operatorname{sech}^2v$. Its integral over $\mathbb R$ is $-2$, and the angular integral contributes $2\pi$. This also shows that total curvature is unchanged by the dilation parameter $a$.
\end{proof}

\begin{proposition}[The complete catenoid is unstable]\label{prop:catenoid-unstable-example}
The catenoid admits a compactly supported smooth normal variation with negative second variation. Each of its open halves $v>0$ and $v<0$, however, is stable for variations compactly supported in that half.
\end{proposition}
\begin{proof}
The periodic unit normal $(\operatorname{sech}v\cos\theta,\operatorname{sech}v\sin\theta,-\tanh v)$ makes these scalar normal variations globally defined. For an angularly constant smooth compactly supported test $f(v)$, the metric and curvature calculations give
\[
 Q(f,f)=2\pi\int_{\mathbb R}\bigl((f')^2-2\operatorname{sech}^2v\,f^2\bigr)\,dv.
\]
Take $f_R$ equal to one on $[-R,R]$, zero outside $[-2R,2R]$, with $|f_R'|\leq C/R$. The positive term is $O(R^{-1})$ and dominated convergence makes the negative integral tend to $-4$, so $Q(f_R,f_R)<0$ for large $R$. Each open half is a minimal graph by Example \ref{ex:explicit-catenoid}; its calibration in its graph cylinder, Proposition \ref{prop:graph-calibration}, gives area minimization for compactly supported variations and hence nonnegative second variation. This proves instability and half-surface stability; it does not compute the full Morse index.
\end{proof}

\begin{proposition}[Rotational minimal hypersurfaces and their first integral]\label{prop:rotational-catenoid-ode}
For $n\geq2$, let $r:I\to(0,\infty)$ be smooth and parametrize a hypersurface in $\mathbb R^{n+1}$ by $X(z,\omega)=(r(z)\omega,z)$, $\omega\in S^{n-1}$. It is minimal if and only if
\[
 rr''=(n-1)(1+(r')^2).
\]
Along any such solution,
\[
 \frac{r^{n-1}}{\sqrt{1+(r')^2}}=c>0
\]
is constant. The constant here is the reciprocal of some source conventions for the flux parameter.
\end{proposition}
\begin{proof}
The area density, after integration over the unit sphere, is $r^{n-1}\sqrt{1+(r')^2}$. Its Euler--Lagrange equation simplifies to the stated ODE. To see that radial stationarity is sufficient, the normal vector is $(\omega,-r')/\sqrt{1+(r')^2}$. A radial variation $\eta(z)\omega$ has nonzero normal factor $\eta/\sqrt{1+(r')^2}$, and the scalar mean curvature is constant on each rotational orbit. Vanishing against all $\eta$ therefore forces the scalar mean curvature to vanish. Finally direct differentiation gives
\[
 \left(\frac{r^{n-1}}{\sqrt{1+(r')^2}}\right)'
 =\frac{r^{n-2}r'}{(1+(r')^2)^{3/2}}
 \bigl((n-1)(1+(r')^2)-rr''\bigr)=0.
\]
This derivation fixes the constant without dividing by $r'$ at the waist. Compare \cite{ChMin}, Section 14, pp.~37--40.
\end{proof}

\begin{example}[Higher-dimensional catenoids fit in a finite-height slab]\label{ex:higher-catenoid-height}
Let $r(0)=a>0$, $r'(0)=0$ in Proposition \ref{prop:rotational-catenoid-ode}. The even solution has upper half described by
\[
 z=\int_a^{r(z)}\frac{d\rho}{\sqrt{(\rho/a)^{2n-2}-1}}.
\]
When $n=2$, its height tends to infinity and $r(z)=a\cosh(z/a)$. When $n\geq3$, the limiting height is finite:
\[
 Z_a=a\int_1^\infty\frac{ds}{\sqrt{s^{2n-2}-1}}<\infty.
\]
In every dimension the resulting hypersurface is complete and proper, with its two ends escaping horizontally; the finite height for $n\geq3$ is not a finite intrinsic distance to an edge.
\end{example}
\begin{proof}
Smooth ODE existence and uniqueness apply near $(r,r')=(a,0)$ since $r>0$. The equation gives $r''>0$, so $r'>0$ for $z>0$, and uniqueness gives evenness. The constant is $c=a^{n-1}$. Solving the first integral for $r'$ and separating variables yields the formula. Near $s=1$ the integrand is bounded by a constant times $(s-1)^{-1/2}$, which is integrable. At infinity it is comparable to $s^{-(n-1)}$, giving divergence for $n=2$ and convergence for $n\geq3$. The solution can continue as long as $r$ stays bounded, because its derivative is then bounded by the first integral. Thus its end has $r\to\infty$. Intrinsic lengths dominate changes in the radial coordinate, so each end has infinite distance. Properness follows from the same radial divergence and compactness of the spherical factor. Compare \cite{ChMin}, Exercises 14.2--14.3.
\end{proof}

\begin{example}[The helicoid: a complete embedded minimal plane]\label{ex:explicit-helicoid}
For $a>0$, define
\[
 X_a(t,\theta)=a(t\cos\theta,t\sin\theta,\theta),\qquad (t,\theta)\in\mathbb R^2.
\]
It is a proper embedded complete minimal surface, with
\[
 g=a^2\bigl(dt^2+(1+t^2)d\theta^2\bigr),\qquad
 K=-\frac1{a^2(1+t^2)^2}.
\]
In particular, unlike the catenoid, it has the topology of a plane and has infinite absolute total curvature.
\end{example}
\begin{proof}
Direct differentiation gives the metric. Substitute $t=\sinh v$; the metric becomes $a^2\cosh^2v(dv^2+d\theta^2)$ and the three coordinate functions become harmonic, proving minimality by Lemma \ref{lem:example-conformal-harmonic}. The conformal Gauss curvature calculation gives the displayed $K$. Height determines $\theta$ and then the horizontal coordinates determine $t$, proving injectivity. A bounded image bounds both parameters, proving properness; the metric dominates $a^2(dt^2+d\theta^2)$, proving completeness. On $|v|<1$ the integral of $|K|dA=\operatorname{sech}^2v\,dv\,d\theta$ is infinite because $\theta$ ranges over all of $\mathbb R$. See \cite{ChMin}, Section 13 and Example 15.6.
\end{proof}

\begin{proposition}[The catenoid-helicoid association and its period obstruction]\label{prop:catenoid-helicoid-associate}
On the parameter plane with $z=\theta+iv$, set
\[
 \Phi(z)=a(\cos z,\sin z,-iz),\qquad
 X_\alpha=\operatorname{Re}(e^{-i\alpha}\Phi),\quad\alpha\in\mathbb R.
\]
Every $X_\alpha$ is a conformal minimal immersion with metric $a^2\cosh^2v(d\theta^2+dv^2)$. The choices $\alpha=0$ and $\alpha=\pi/2$ give a catenoid covering and a helicoid, up to a rigid motion. The immersion descends to the cylinder $\theta\pmod{2\pi}$ if and only if $\sin\alpha=0$.
\end{proposition}
\begin{proof}
The holomorphic derivative satisfies $\Phi'\cdot\Phi'=a^2(\sin^2z+\cos^2z-1)=0$ and $|\Phi'|^2=2a^2\cosh^2v$. Taking real and imaginary parts therefore gives orthogonal real derivatives of equal squared norm $|\Phi'|^2/2$. Multiplication by a unit complex number preserves both identities and harmonicity. The imaginary part is $a(-\sin\theta\sinh v,\cos\theta\sinh v,-\theta)$, a rotated/reflected helicoid. Finally $\Phi(z+2\pi)-\Phi(z)=(0,0,-2\pi ia)$, whose real part after multiplication by $e^{-i\alpha}$ is $(0,0,-2\pi a\sin\alpha)$. This vanishes exactly as claimed. Local isometry does not identify the global topologies. Compare \cite{ChMin}, Lemma 15.5 and Example 15.6.
\end{proof}

\begin{example}[Enneper's surface is immersed and complete but not embedded]\label{ex:explicit-enneper}
The polynomial map
\[
 E(u,v)=\left(u-\frac{u^3}{3}+uv^2,
 -v+\frac{v^3}{3}-u^2v,u^2-v^2\right)
\]
is a proper complete conformal minimal immersion of $\mathbb R^2$ into $\mathbb R^3$, with metric $(1+u^2+v^2)^2(du^2+dv^2)$. It is not injective: $E(\sqrt3,0)=E(-\sqrt3,0)=(0,0,3)$.
\end{example}
\begin{proof}
Differentiation gives $|E_u|^2=|E_v|^2=(1+u^2+v^2)^2$ and $\langle E_u,E_v\rangle=0$; each coordinate is harmonic. The positive metric factor proves immersion and Lemma \ref{lem:example-conformal-harmonic} proves minimality. The metric dominates the Euclidean one, giving completeness. With $z=u+iv$, the first two components satisfy $E_1+iE_2=\overline z-z^3/3$, hence $|E_1+iE_2|\geq |z|^3/3-|z|$. This tends to infinity, proving properness. The two displayed parameters prove failure of injectivity. At them, $E_u=(-2,0,\pm2\sqrt3)$ and $E_v=(0,-4,0)$ span different tangent planes, so the image itself has a crossing and is not an embedded surface. This is Enneper's example from \cite{ChMin}, Section 15.5, in a fixed polynomial normalization.
\end{proof}

\begin{proposition}[Total curvature counts the parameter surface]\label{prop:enneper-total-curvature}
For the Enneper immersion,
\[
 K=-\frac4{(1+u^2+v^2)^4},\qquad
 \int_{\mathbb R^2}K\,dA=-4\pi.
\]
The integral uses the induced metric on the parameter plane; self-intersections do not cancel area or curvature contributions.
\end{proposition}
\begin{proof}
Put $r^2=u^2+v^2$ and $\lambda=1+r^2$. The identity $\Delta\log(1+r^2)=4/(1+r^2)^2$ and the conformal curvature formula give $K$. Multiplying by $dA=(1+r^2)^2du\,dv$ and integrating in polar coordinates yields $-8\pi\int_0^\infty r(1+r^2)^{-2}\,dr=-4\pi$. Integration over the immersion always retains parameter multiplicity.
\end{proof}

\begin{remark}[Other explicit models already developed in the project]\label{rem:example-minimal-links}
Example \ref{ex:scherk-graph} computes the minimal equation for Scherk's graph and its loss of a global gradient bound. Example \ref{ex:complex-cusp-current} constructs a higher-codimension calibrated minimizer, and Propositions \ref{prop:cusp-density} and \ref{prop:cusp-tangent} compute its density and double-plane tangent. These entries retain their original tags and locations. The present catenoid calculation shows explicitly why the words minimal, stable, and area minimizing must be distinguished.
\end{remark}
